\section*{Минимальный двудольный \(C_4\)-лифт}

Полный перебор помеченных четырёхвершинных графов показывает, что квадрат
\(0-1-2-3-0\) является единственным минимальным двудольным циклом,
содержащим цепь \(A_3\). Все его рёбра могут быть нечётными относительно
\(\Gamma_4=\operatorname{diag}(1,-1,1,-1)\), а Schur-комплемент координатной
вершины порождает эффективную комплексную хорду.

Однако канонические три поколения образуют не координатные вершины
\(0,1,2\), а sum-zero подпространство, ортогональное
\(u=(1,1,1,1)^T/2\). Для квадрата с потоком
\[
 \left\|(I-uu^\dagger)A_{C_4}(\Phi)u\right\|^2
 =\frac{(1-\cos\Phi)(3+\cos\Phi)}4.
\]
Ненулевой поток смешивает affine-синглет и триплет, а
\(\|[\Gamma_4,uu^\dagger]\|_F^2=2\). Кроме того, исключаемая координатная
вершина содержит только четверть синглетной нормы. Поэтому координатный
Schur-комплемент не является редукцией канонического механизма \(1+3\).

Квадрат чинит двудольность графа, но не сохраняет прежний вывод числа
поколений. Ветвь требует либо большего градуированного пространства, либо
полного повторного вывода поколений на четырёхмерном носителе.